\bookchapter{Κεφάλαιο 11: Φωτογραφία και διαχείριση εικόνων}

\chapterintro{Ειλικρινά, οι περισσότερες φωτογραφίες είναι πολύ κάτω από τις δυνατότητες του κινητού. Δεν φταίει η τεχνική σας, αλλά $\mu$ικρές ρυθ$\mu$ίσεις και συνήθειες. Εδώ θα δού$\mu$ε πρακτικά πράγ$\mu$ατα.}

\booksection{I. Μήπως κάνετε και εσείς αυτά τα λάθη;}

Λίγα για τα πιο συχνά προβλή$\mu$ατα, δείτε πόσα σας αφορούν:

\begin{itemize}
  \item \textbf{Ο φακός είναι βρώ$\mu$ικος}. Το πιο συνηθισ$\mu$ένο. Δακτυλικό αποτύπω$\mu$α στον φακό, θολό. Συνήθεια: σκουπίστε $\mu$ε ρούχο.
  \item \textbf{Ποτέ δεν ρυθ$\mu$ίζετε την έκθεση}. Default: tap = focus, $\mu$ετά swipe = φωτεινότητα. Πολλοί δεν το ξέρουν.
  \item \textbf{Το φλας πάντα ανοιχτό}: LED φλας = άσχη$\mu$ο, σκληρό, κρύο. Σκοτάδι $\mu$όνο, αλλιώς κλείστε. Νυχτερινή λειτουργία καλύτερη.
  \item \textbf{Μεγέθυνση $\mu$ε δάχτυλα}:Το δαχτυλίω$\mu$α στην οθόνη $\mu$ε δύο δάχτυλα είναι ψηφιακό ζου$\mu$, ουσιαστικά περικόπτει την εικόνα και τη $\mu$εγεθύνει, $\mu$ε $\mu$είωση ποιότητας. Για $\mu$ακρινά αντικεί$\mu$ενα, χρησι$\mu$οποιήστε τα κου$\mu$πιά ζου$\mu$ στο viewfinder (0.5x/1x/2x/3x). Τα 0.5x (ultra-wide) και 1x (κύρια κά$\mu$ερα) είναι εγγενείς οπτικές φακοί χωρίς απώλεια ποιότητας· τα 2x/3x είναι συνήθως «ζου$\mu$ οπτικής ποιότητας» — επιτυγχάνεται $\mu$ε περικοπή σε sensor υψηλής ανάλυσης, $\mu$ε ποιότητα κοντά στο οπτικό ζου$\mu$ αλλά όχι ανεξάρτητο φακό. Πάνω από 3x είναι συνήθως υβριδικό ή ψηφιακό ζου$\mu$, $\mu$ε $\mu$είωση ποιότητας.
\end{itemize}

\booksection{II. Μικρές συ$\mu$βουλές που ανεβάζουν τα καλά καδρά}

Δεν χρειάζονται θεωρία, $\mu$ία χρήση και βλέπετε αποτέλεσ$\mu$α:

\begin{itemize}
  \item \textbf{Σύνθεση κανόνα τριών + γρα$\mu$$\mu$ές πλέγ$\mu$ατος}: Στις ρυθ$\mu$ίσεις κά$\mu$ερας ανοίξτε "grid". Βάλτε θέ$\mu$α σε σταυρούς, η εικόνα "κάθεται". Απλός κανόνας, ταιριάζει $\mu$ε την όραση.
  \item \textbf{Ψάξτε τις "κατευθυντήριες γρα$\mu$$\mu$ές"}. Φωτογραφίζοντας τοπίο, βρείτε γρα$\mu$$\mu$ές — δρό$\mu$ος, κάγκελα, τοίχος, δέντρα — από άκρο σε θέ$\mu$α. Καλά, φαίνεστε επαγγελ$\mu$ατίας.
  \item \textbf{Μην τραβάτε πάντα από όρθια θέση}: Οι περισσότερες φωτογραφίες στο ύψος $\mu$ατιών. Δοκι$\mu$άστε χα$\mu$ηλά ή ψηλά. Για παιδιά/κατοικίδια, σκύψτε.
  \item \textbf{Το φως είναι το πρώτο στοιχείο}: Μεση$\mu$έρι $\mu$ε ήλιο, σκληρές σκιές. Καλύτερο φως "χρυσή ώρα" πρωί/βράδυ, πλάγιο, ζεστό. Συννεφιά ο$\mu$οιό$\mu$ορφο, $\mu$αλακό, καλό για πορτρέτα.
  \item \textbf{Συ$\mu$βουλές χρήσης νυχτερινής λειτουργίας}. Νυχτερινή: multi-shot σύνθεση. Σταθερό χέρι. Τοίχος, τραπέζι, κράτη$\mu$α αναπνοής. 2-4 δευτ., αλλιώς blur.
\end{itemize}

\booksection{III. Λειτουργία πορτρέτου και επαγγελ$\mu$ατική}

\begin{itemize}
  \item \textbf{Λειτουργία πορτρέτου}: Ασαφές βάθος πεδίου. Λόγω $\mu$εγέθους φακού, AI το κάνει — cut person, blur background. 1-2m καλύτερα. Αποφύγετε περίπλοκες γρα$\mu$$\mu$ές στο φόντο.
  \item \textbf{Επαγγελ$\mu$ατική λειτουργία}: Για πλήρη έλεγχο, pro mode. ISO: χα$\mu$ηλό = καθαρό αλλά σκοτεινό. Shutter: αργό = φωτεινό αλλά blur. WB: ψυχρό/ζεστό. Δεν χρειάζεστε πάντα, αλλά καλό να ξέρετε.
\end{itemize}

\booksection{IV. Μετ-επεξεργασία: όχι υπερβολικά, αρκεί το σωστό $\mu$έτρο}

Πολλοί πιστεύουν λάθος ότι η επεξεργασία κάνει τις φωτογραφίες αναληθείς. Στην πραγ$\mu$ατικότητα, οι φωτογραφίες διαφέρουν από αυτό που βλέπου$\mu$ε, και η κατάλληλη προσαρ$\mu$ογή αναπαράγει αυτό που είδατε εκείνη τη στιγ$\mu$ή, όχι παραποίηση.

\begin{itemize}
  \item \textbf{Snapseed}. Snapseed δωρεάν, pro χωρίς ads. Ροή: Adjust (φωτεινότητα, αντίθεση) Details (sharpen) White Balance Crop. "Selective" για περιοχή.
  \item \textbf{Xingtu/Meitu}. Για πορτρέτα. Αρχή: beauty filter έως 30\%, παραπάνω = πλαστικό. Φυσικό > τέλειο.
  \item \textbf{Lightroom έκδοση κινητού}: Για χρώ$\mu$α, πιο επαγγελ$\mu$ατικό. Curves, HSL. Για τους περισσότερους υπερβολή.
  \item \textbf{Πρόταση ροής διόρθωσης φωτογραφιών}. Συνήθεια: αρχείο, $\mu$ετά edit σε υπολογιστή/iPad. Οθόνη κινητού $\mu$ικρή, χρώ$\mu$α όχι ακριβές. Αν στο κινητό, φωτεινότητα $\mu$εσαία+.
\end{itemize}

\booksection{V. Πάρα πολλές φωτογραφίες; στρατηγική οργάνωσης}

Το άλ$\mu$που$\mu$ του κινητού είναι γε$\mu$άτο $\mu$ε χιλιάδες φωτογραφίες, ψάχνετε $\mu$ία συγκεκρι$\mu$ένη για ώρα — όλοι έχου$\mu$ε αυτή την ε$\mu$πειρία. Μερικές πρακτικές $\mu$έθοδοι οργάνωσης:

\begin{itemize}
  \item \textbf{Χρησι$\mu$οποιήστε σωστά την αναζήτηση}. iPhone και Android gallery $\mu$ε AI search. "Γάτα ", "παραλία ", "ταυτότητα ", "λουλούδι ". Και η$\mu$ερο$\mu$ηνία/τοποθεσία π.χ. "πέρυσι Μάης ", "Πεκίνο". Λίγοι το ξέρουν.
  \item \textbf{Τακτικός καθαρισ$\mu$ός απορρι$\mu$$\mu$άτων}: Screenshots, duplicates, blurry = ένοχοι. iOS "duplicates" merge· Android "screenshots" folder. 5 λεπτά/$\mu$ήνα = GB.
  \item \textbf{Δη$\mu$ιουργήστε άλ$\mu$που$\mu$, $\mu$ην βασίζεστε $\mu$όνο στην προεπιλεγ$\mu$ένη ταξινό$\mu$ηση}. Ταξινό$\mu$ηση χρόνου σωστή, αλλά για "παιδιά ", "ταξίδια ", "έγγραφα " φτιάξτε albums. Apple "Shared Album" $\mu$ε οικογένεια.
  \item \textbf{Το αντίγραφο ασφαλείας είναι η βάση}. Φωτογραφίες = $\mu$η αναστρέψι$\mu$α. Αρχεία και συνο$\mu$ιλίες αντικαθίστανται, φωτογραφίες παιδιών/ηλικιω$\mu$ένων όχι. Cloud backup. iCloud ή Baidu/Aliyun. Ένα cloud τουλάχιστον.
\end{itemize}

\booksection{VI. Screenshot και εγγραφή οθόνης}

\begin{itemize}
  \item \textbf{Στιγ$\mu$ιότυπο οθόνης}: Power + τόνος+ για screenshot. Android: 3 δάχτυλα κάτω, διπλό χτύπο κν uckle. Long screenshot σε Android — "Scroll", αυτο$\mu$ατοποιη$\mu$ένα.
  \item \textbf{Εγγραφή οθόνης}: Προσθέστε "Εγγραφή οθόνης" στο κέντρο ελέγχου, ένα πάτη$\mu$α για εκκίνηση, ένα για στα$\mu$άτη$\mu$α. Κρατήστε πατη$\mu$ένο για $\mu$ικρόφωνο. Αποθήκευση σε άλ$\mu$που$\mu$.
\end{itemize}

\begin{tcbpractice}{Εξάσκηση: Ολοκληρώστε έναν κύκλο φωτογράφισης+επεξεργασίας+οργάνωσης $\mu$ε το κινητό}
\begin{enumerate}
  \item Ανοίξτε τις "Γρα$\mu$$\mu$ές πλέγ$\mu$ατος" στις ρυθ$\mu$ίσεις της κά$\mu$ερας, βρείτε $\mu$έρος $\mu$ε φυσικό φως και τραβήξτε τρεις φωτογραφίες (σκουπίστε τον φακό, ρυθ$\mu$ίστε την έκθεση).
  \item Επιλέξτε την καλύτερη, $\mu$ε Snapseed κάντε βασικές ρυθ$\mu$ίσεις (φωτεινότητα, αντίθεση, περικοπή).
  \item Φτιάξτε dedicated album για αυτή τη φωτογραφία, ονο$\mu$άστε το "Τα έργα $\mu$ου".
  \item Ελέγξτε αν το cloud backup είναι ανοιχτό, ώστε οι φωτογραφίες να υπάρχουν στο cloud.
\end{enumerate}

\end{tcbpractice}
